\section{Hardware}
\label{sec:hardware}

Everything in the earlier sections --- the magnitude, the Welch spectrum, the
freeze index, the gate-and-debounce state machine --- has to run on something small
enough to live on an ankle and cheap enough to build in a one-week sprint. This
section describes that physical build: the single sensing board, the one external
display, the optional wireless bridge, how it all sits inside an ankle sleeve, and
the two firmware ``personalities'' the same board can wear.

\begin{keyidea}
Cadence is a \textbf{monitor}, not a therapy. A coin vibration motor is physically
fitted to the build, but it is \textbf{disabled in firmware}: the device never buzzes,
cues, or stimulates. Every freeze it finds is \emph{reported} --- on an OLED, on the
onboard LEDs, and on a live web dashboard --- and nothing more. Read every mention of
the motor below as ``present on the board, held off in software.''
\end{keyidea}

\subsection{The one board that does the work: Circuit Playground Express}

The heart of Cadence is a single \textbf{Adafruit Circuit Playground Express (CPX)}.
It is a 51\,mm disc that already carries everything the core detector needs, so the
sensing, the maths, and the local display all happen on one part:

\begin{itemize}
  \item a \textbf{SAMD21} Arm Cortex-M0+ microcontroller --- the CPU that runs the DSP;
  \item an onboard \textbf{LIS3DH} 3-axis MEMS accelerometer (I\textsuperscript{2}C),
        the only sensor Cadence uses, sampled at \textbf{64\,Hz} on the
        \(\pm 8\,\mathrm{g}\) range;
  \item a ring of \textbf{10 NeoPixels} that we drive as a colour-coded state light;
  \item two push-buttons and a \textbf{slide switch} we repurpose as mode/enable inputs;
  \item 2\,MB of SPI flash for untethered logging.
\end{itemize}

\begin{intuition}
Why one disc instead of a pile of breakout boards? In a one-week wearable sprint the
enemy is wiring. The CPX folds the MCU, the accelerometer, and the status LEDs into a
part you can sew flat against the limb with \emph{zero} jumper wires for the core job.
The only wires we add are the four that drive an external screen --- and even those
are optional.
\end{intuition}

The \(\pm 8\,\mathrm{g}\) range is a deliberate choice, not a default. Heel-strike
transients at the ankle routinely exceed \(\pm 2\,\mathrm{g}\), so a \(\pm 2\,\mathrm{g}\)
setting would clip the very peaks that separate vigorous gait from a freeze. The
\(64\,\mathrm{Hz}\) sample rate gives a Nyquist limit of \(32\,\mathrm{Hz}\), which sits
comfortably above the highest band of interest (the \(3\text{--}8\,\mathrm{Hz}\) freeze
band, and the \(4\text{--}6\,\mathrm{Hz}\) rest-tremor band), so nothing of clinical
interest aliases.

\subsection{Pins and connections}

The wiring is dominated by one addition: an external \textbf{SSD1306} OLED for the
live freeze read-out. The CPX's onboard I\textsuperscript{2}C bus is already busy with
the LIS3DH accelerometer, so the OLED is driven on a \emph{second} I\textsuperscript{2}C
bus brought out on the CPX pads --- \textbf{SDA = A5}, \textbf{SCL = A4} --- using a
separate \code{Wire} instance. Keeping the screen off the sensor's bus means a slow or
chatty display can never stall an accelerometer read. Figure~\ref{fig:wiring} gives the
full pad-by-pad map of the build.

\begin{figure}[H]
  \centering
  \includegraphics[width=0.86\linewidth]{wiring_pinout}
  \caption{Pad-by-pad wiring of the Cadence build. The CPX senses, runs the freeze-index
  detector, drives the onboard NeoPixel ring, and brings out a second
  I\textsuperscript{2}C bus (A5/A4) for the external SSD1306 OLED; three UART wires
  (A7/A6 + a shared ground, 3.3\,V logic, no level shifter) form the optional link to an
  ESP32 Wi-Fi bridge. \emph{Note:} this sheet is inherited from an earlier closed-loop
  iteration and still annotates the coin motor as a ``cue'' on A1; in shipped Cadence
  firmware that motor pin is held off --- the device is a monitor.}
  \label{fig:wiring}
\end{figure}

\begin{table}[H]
  \centering
  \small
  \caption{Key connections on the Circuit Playground Express. The first block is the
  always-present core; the second block is the optional wireless bridge.}
  \label{tab:pins}
  \begin{tabular}{@{}llp{6.2cm}@{}}
    \toprule
    \textbf{CPX pin / part} & \textbf{Direction} & \textbf{Role in Cadence} \\
    \midrule
    LIS3DH accel  & onboard I\textsuperscript{2}C & 3-axis sensor, 64\,Hz, \(\pm 8\,\mathrm{g}\) --- the only sensor used \\
    NeoPixels \(\times 10\) & onboard out & state-colour ring (still / walking / freeze) \\
    Slide switch  & onboard in & selects RECORDER vs.\ MONITOR role \\
    Button A      & onboard in & hold-still calibration of the still-floor gate \\
    A5 (SDA2)     & out (I\textsuperscript{2}C) & external SSD1306 OLED data \\
    A4 (SCL2)     & out (I\textsuperscript{2}C) & external SSD1306 OLED clock \\
    A1            & out (held off) & coin vibration motor --- \textbf{disabled in firmware} \\
    GND           & power (0\,V) & common ground (shared with the ESP32) \\
    \addlinespace
    A7            & UART TX & \(a_x,a_y,a_z\) stream out to ESP32 (optional bridge) \\
    A6            & UART RX & control bytes in from ESP32 (optional bridge) \\
    micro-USB     & 5\,V in & power; also the USB link for log download \\
    \bottomrule
  \end{tabular}
\end{table}

\subsection{Telling the wearer: NeoPixels and the OLED}

Cadence has two local outputs, both read-only indicators of the same state machine:

\begin{itemize}
  \item the \textbf{NeoPixel ring} is a glanceable colour light --- it shows the
        headline \textsc{still} / \textsc{walking} / \textsc{freeze} verdict at arm's
        length, with red reserved for a committed freeze;
  \item the \textbf{SSD1306 OLED} carries the detail the LEDs cannot: the live freeze
        index, the movement-energy gate state, and a text label of the current phase.
\end{itemize}

Both are downstream of the gate \(\rightarrow\) FI \(\rightarrow\) debounce pipeline of
Section~\ref{sec:detector}; neither feeds back into the decision. Putting the screen on
its own bus (A5/A4) is what lets the OLED be as slow as it likes without ever delaying a
64\,Hz accelerometer sample.

\subsection{Going wireless: the optional 3-board topology}

The standalone build is fully untethered: the CPX senses, decides, and displays on the
ankle, powered from a small battery over USB. For a richer, screen-rich \emph{live}
view --- the web dashboard of Section~\ref{sec:dashboard} --- Cadence offers an optional
three-board topology that adds an \textbf{ESP32} as a Wi-Fi bridge between the ankle and
a laptop.

\begin{figure}[H]
  \centering
  \resizebox{0.96\textwidth}{!}{%
  \begin{tikzpicture}[
      font=\sffamily,
      box/.style={rounded corners=3pt,draw=#1,line width=1pt,fill=#1!8,
                  minimum width=3.5cm,minimum height=1.9cm,align=center,inner sep=5pt},
      lbl/.style={font=\footnotesize\bfseries},
      sub/.style={font=\scriptsize\color{mute}},
    ]
    % --- the three boards, generously spaced ---
    \node[box=teal]  (cpx)    {\textbf{CPX (ankle)}\\[2pt]
                               {\scriptsize\color{mute} LIS3DH @ 64\,Hz}\\
                               {\scriptsize\color{mute} on-board FI detector}\\
                               {\scriptsize\color{mute} NeoPixels + OLED}};
    \node[box=navy,right=3.0cm of cpx] (esp) {\textbf{ESP32}\\[2pt]
                               {\scriptsize\color{mute} UART\,$\leftrightarrow$\,Wi-Fi bridge}\\
                               {\scriptsize\color{mute} relays the stream}\\
                               {\scriptsize\color{mute} no on-board decision}};
    \node[box=steel,right=3.0cm of esp] (host) {\textbf{Laptop}\\[2pt]
                               {\scriptsize\color{mute} dashboard server}\\
                               {\scriptsize\color{mute} optional CNN}\\
                               {\scriptsize\color{mute} renders live panels}};

    % --- forward data path: accel stream travels right ---
    \draw[-{Latex[length=2.6mm]},teal,line width=1.1pt]
      ($(cpx.east)+(0,0.32)$) -- node[lbl,above,yshift=1pt,teal]{UART}
      node[sub,below,yshift=-1pt]{$a_x,a_y,a_z$ @ 64\,Hz}
      ($(esp.west)+(0,0.32)$);
    \draw[-{Latex[length=2.6mm]},navy,line width=1.1pt]
      ($(esp.east)+(0,0.32)$) -- node[lbl,above,yshift=1pt,navy]{Wi-Fi / TCP}
      node[sub,below,yshift=-1pt]{packets}
      ($(host.west)+(0,0.32)$);

    % --- return control path: thin, runs back left ---
    \draw[-{Latex[length=2.4mm]},freeze,densely dashed,line width=0.9pt]
      ($(host.west)+(0,-0.42)$) --
      node[sub,below,yshift=-1pt,freeze]{control / status} ($(esp.east)+(0,-0.42)$);
    \draw[-{Latex[length=2.4mm]},freeze,densely dashed,line width=0.9pt]
      ($(esp.west)+(0,-0.42)$) -- ($(cpx.east)+(0,-0.42)$);

    % --- annotations under the row ---
    \node[sub,below=0.7cm of cpx.south,align=center]
      {sensing + decision\\ on the ankle};
    \node[sub,below=0.7cm of esp.south,align=center]
      {dumb relay\\ (3.3\,V logic,\\ shared ground)};
    \node[sub,below=0.7cm of host.south,align=center]
      {display + analytics\\ off the body};
  \end{tikzpicture}}
  \caption{The optional 3-board topology: \textbf{CPX}
  \(\;\xrightarrow{\;\text{UART}\;}\;\) \textbf{ESP32}
  \(\;\xrightarrow{\;\text{Wi-Fi/TCP}\;}\;\) \textbf{laptop}. The accelerometer
  stream flows right (solid); thin control/status bytes return left (dashed). The ESP32
  is a pure bridge --- it makes no decisions. Because both the CPX and ESP32 run at
  3.3\,V logic and share a ground, the UART wires connect directly with no level shifter.}
  \label{fig:topology}
\end{figure}

\begin{intuition}
The split is ``decide near the sensor, display far from it.'' The freeze decision is
small and latency-sensitive, so it stays on the ankle. Pixels, charts, and any
heavier model are big and not safety-critical, so they live on a laptop the wearer
need not carry. The ESP32 in the middle is deliberately dumb: it shuttles bytes and
nothing else, which keeps the part that matters --- the detector --- on one well-tested
board.
\end{intuition}

\subsection{Wearing it: the ankle sleeve}

Cadence is built into a snug compression ankle sleeve with flat internal pockets, so
nothing dangles and the sensor is firmly coupled to the limb. Firm coupling matters:
a loose sensor would swing like a pendulum and pollute the very band the freeze index
reads. The CPX and the battery sit in opposite flat pockets so their masses balance
near the limb axis, the optional motor sits in a thin pocket against the skin, and the
whole assembly is light. Figure~\ref{fig:sleeve} shows the layout from the side, in
cross-section, and as a flat-pocket layer stack.

\begin{figure}[H]
  \centering
  \includegraphics[width=0.97\linewidth]{sleeve_layout}
  \caption{Ankle-sleeve integration: (A) lateral view on the right leg, (B)
  cross-section looking toward the foot, (C) the flat-pocket layer stack. The sleeve is
  snug so the accelerometer reads true limb motion rather than pendulum swing; CPX and
  battery sit on opposite faces to balance; the build sits flush with nothing to snag.
  The vibration motor shown against the skin is fitted but, in Cadence's monitor
  firmware, held off.}
  \label{fig:sleeve}
\end{figure}

Power is a small LiPo cell delivered to the CPX over its USB connector; the same USB
connection doubles as the link for pulling logged data off the board. In the wireless
build the ESP32 is powered alongside the CPX and shares its ground, which is exactly why
the UART lines need no level translation.

\subsection{Two firmware roles on one board}

The same CPX wears one of two firmware ``personalities,'' chosen so the project can both
\emph{gather} labelled data and \emph{run live} without changing hardware. The slide
switch selects between them.

\begin{table}[H]
  \centering
  \small
  \caption{The two firmware roles. One board, one accelerometer --- two jobs.}
  \label{tab:roles}
  \begin{tabular}{@{}p{2.4cm}p{6.0cm}p{6.0cm}@{}}
    \toprule
     & \textbf{RECORDER} & \textbf{MONITOR} \\
    \midrule
    Sketch    & \code{cpx\_fog\_logger.ino} & \code{cpx\_fog\_standalone.ino} \\
    Purpose   & capture raw accel + auto-labelled phases for analysis & live freeze detection and display \\
    Trigger   & flip the slide switch to start/stop logging & runs continuously on power-up \\
    Stores to & 2\,MB onboard SPI flash, untethered & nothing --- decides every 2\,s in RAM \\
    Output    & a file pulled later over USB with \code{cpx\_dump.py} & live OLED text + NeoPixel state colours \\
    Motor     & n/a & fitted but disabled \\
    \bottomrule
  \end{tabular}
\end{table}

The \textbf{RECORDER} role is how the project earns its right to claim real numbers: a
clinician or wearer flips the slide switch, walks (and freezes), and the board logs
accelerometer samples with auto-labelled phases to flash with no laptop attached. Later
the file is pulled over USB with \code{cpx\_dump.py} for the offline analysis behind every
result in this report. The \textbf{MONITOR} role is the shipped device: it runs the
gate~\(\rightarrow\)~FI~\(\rightarrow\)~debounce pipeline live and paints the verdict onto
the OLED and LEDs.

\subsection{Does it fit on the chip?}

A fair question for a Cortex-M0+ with no floating-point unit: can the SAMD21 finish a
window of DSP inside its budget? The arithmetic is comfortable. A decision is due only
once every \textbf{2\,s} (the 128-sample hop), and the work per decision is one
Hann-tapered Welch PSD over a 256-sample window plus a handful of band-power sums and a
single ratio. That is a few thousand multiply-adds spread across a two-second window ---
well within the SAMD21's headroom, leaving the CPU mostly idle between decisions to
service the 64\,Hz sampling, the LEDs, and the OLED. Figure~\ref{fig:flow} traces that
whole MONITOR loop end to end.

\begin{keyidea}
The cadence of the whole device is set by the 2\,s hop, not by raw compute. The long
4\,s window is what buys the spectral resolution to separate the \(3\text{--}8\,\mathrm{Hz}\)
freeze band from the walking rhythm; the 2\,s hop is how often that window is re-scored.
Both numbers are about \emph{seeing the freeze clearly}, and both leave the little SAMD21
with time to spare.
\end{keyidea}

\begin{figure}[H]
  \centering
  \includegraphics[height=0.73\textheight]{coding_flowchart}
  \caption{The MONITOR firmware loop (\code{cpx\_fog\_standalone.ino}): read the three
  axes at 64\,Hz, form the magnitude, keep the last 256 samples (4\,s), and once every
  2\,s compute the freeze index and movement energy, apply the still-floor gate, the
  \(\FI>1.815\) test and the assert/release-after-two debounce, then update the
  NeoPixels and log line. \emph{Note:} this diagram is inherited from a closed-loop
  iteration and still draws a ``motor cue'' block; in shipped Cadence that output path
  is held off --- only the NeoPixel/OLED/serial update is active.}
  \label{fig:flow}
\end{figure}

\begin{takeaway}
One 51\,mm board does the sensing and the decision; one external OLED and the onboard
LEDs report it; an optional ESP32 carries it to a laptop dashboard. The coin motor is
on the board but off in code --- Cadence watches for freezes and tells you, and that is
all.
\end{takeaway}
